\documentclass[12pt, a4paper]{report}

% --- Typography & Encoding ---
\usepackage[T1]{fontenc}
\usepackage[utf8]{inputenc}
\usepackage{newtxtext,newtxmath} % Times New Roman-like typography
\usepackage{microtype}
\usepackage{setspace}
\onehalfspacing

% --- Page Geometry ---
\usepackage[
  a4paper,
  left=1.00in,
  right=0.85in,
  top=0.78in,
  bottom=0.80in,
  headheight=16pt,
  headsep=8pt,
  footskip=28pt
]{geometry}

% --- Custom Chapter/Section Formats ---
\usepackage{titlesec}
\titleformat{\chapter}[hang]
  {\bfseries\fontsize{16}{19}\selectfont}
  {\thechapter.}{0.55em}{\MakeUppercase}
\titlespacing*{\chapter}{0pt}{0pt}{20pt}

\titleformat{\section}
  {\bfseries\fontsize{13}{16}\selectfont}
  {\thesection.}{0.55em}{}
\titlespacing*{\section}{0pt}{16pt}{8pt}

\titleformat{\subsection}
  {\bfseries\fontsize{12}{15}\selectfont}
  {\thesubsection.}{0.55em}{}
\titlespacing*{\subsection}{0pt}{14pt}{6pt}

% Simplify chapter headers to remove "Chapter" prefix
\makeatletter
\renewcommand{\@makechapterhead}[1]{%
  \vspace*{0pt}%
  {\parindent\z@ \raggedright \normalfont
    \bfseries\fontsize{16}{19}\selectfont
    \thechapter.\quad\MakeUppercase{#1}\par\nobreak
    \vskip 20pt}}
\makeatother

% --- Page Borders (MACE Style) ---
\usepackage{tikz}
\usepackage{eso-pic}
\AddToShipoutPictureBG{%
\begin{tikzpicture}[remember picture,overlay]
\draw[line width=1.0pt]
    ([xshift=0.3in,yshift=-0.3in]current page.north west)
    rectangle
    ([xshift=-0.3in,yshift=0.3in]current page.south east);
\end{tikzpicture}%
}

% --- Header and Footer (User Request) ---
\usepackage{fancyhdr}
\newcommand{\ReportTitle}{ClauseIQ -- Machine Learning-Based Contract Clause Classification and Intelligent Search System}
\newcommand{\DepartmentName}{Department of Computer Applications, MACE}

\pagestyle{fancy}
\fancyhf{}
\fancyhead[R]{\small\ReportTitle}
\fancyfoot[L]{\small\DepartmentName}
\fancyfoot[R]{\small\thepage}
\renewcommand{\headrulewidth}{0.6pt}
\renewcommand{\footrulewidth}{0.6pt}

% Apply same style on chapter-opening pages
\fancypagestyle{plain}{%
  \fancyhf{}
  \fancyhead[R]{\small\ReportTitle}
  \fancyfoot[L]{\small\DepartmentName}
  \fancyfoot[R]{\small\thepage}
  \renewcommand{\headrulewidth}{0.6pt}
  \renewcommand{\footrulewidth}{0.6pt}
}

% --- Paragraph Formatting ---
\setlength{\parindent}{0pt}
\setlength{\parskip}{0.75em}
\usepackage{ragged2e}
\AtBeginDocument{\justifying}

% --- Code Snippet Styling ---
\usepackage{xcolor}
\usepackage{listings}
\definecolor{codebg}{gray}{0.14}
\definecolor{codetext}{gray}{0.96}
\lstset{
  basicstyle=\ttfamily\small\color{codetext},
  backgroundcolor=\color{codebg},
  frame=single,
  rulecolor=\color{codebg},
  breaklines=true,
  columns=fullflexible,
  keepspaces=true,
  showstringspaces=false,
  xleftmargin=0.04\textwidth,
  xrightmargin=0.04\textwidth,
  aboveskip=8pt,
  belowskip=4pt
}

% --- Mathematics, Lists, Tables ---
\usepackage{amsmath, amssymb}
\usepackage{booktabs}
\usepackage{array}
\usepackage{longtable}
\usepackage{tabularx}
\usepackage{enumitem}
\setlist[itemize]{leftmargin=1.5em,itemsep=2pt,topsep=4pt}
\setlist[enumerate]{leftmargin=1.7em,itemsep=2pt,topsep=4pt}

% --- Hyperlinks & References ---
\usepackage[hidelinks]{hyperref}
\usepackage{xurl}

\begin{document}

% ==========================================
% CHAPTER 1: INTRODUCTION
% ==========================================
\chapter{INTRODUCTION}
\label{ch:introduction}

\section{Contract Review Problems in Industries}
Commercial contracts represent the legal and operational backbone of virtually all corporate business transactions. These documents, which include Non-Disclosure Agreements (NDAs), Service Level Agreements (SLAs), Master Service Agreements (MSAs), joint venture structures, and licensing arrangements, outline the rights, duties, and liabilities of the contracting parties. However, despite their critical importance, the administration and evaluation of these documents remain heavily dependent on manual processes in corporate environments. 

Organizations routinely manage thousands of active agreements. As transaction volumes grow, purely manual review does not scale, creating major challenges:
\begin{itemize}
    \item \textbf{Cognitive Fatigue and Errors:} Legal documents are dense, lengthy, and written in highly specialized legal prose. Manual review of hundreds of pages of text causes cognitive fatigue among lawyers, leading to oversight of critical clauses, financial obligations, or non-compliance risks.
    \item \textbf{Inconsistent Interpretations:} Different manual reviewers may interpret identical or similar clauses differently, leading to inconsistent compliance metrics across organization boundaries.
    \item \textbf{High Operational Cost and Bottlenecks:} Due diligence audits and contract reviews require hours of expensive legal labor. The resulting bottlenecks can delay sales cycles and procurement processes.
\end{itemize}

\section{Why Manual Clause Searching is Difficult}
When legal teams or business executives need to locate specific provisions (such as indemnity limits, termination notice periods, or governing law clauses) across historical documents, they face significant hurdles. Manual searching is highly challenging for several key reasons:
\begin{itemize}
    \item \textbf{Lack of Semantic Capabilities:} Standard keyword searches (e.g., using Ctrl+F) look for exact string matches. Legal concepts, however, can be drafted using various synonyms. For example, a search for ``Liability Limit'' will miss clauses titled ``Limitation of Liability'' or sentences referring to ``indemnification caps'' or ``maximum aggregate exposure.''
    \item \textbf{Structural Variation:} Contracts are unstructured text files without a standard schema. Some documents embed key clauses within sections on general covenants, while others isolate them in dedicated appendices.
    \item \textbf{Document Formats:} Legacy contracts are frequently stored as scanned PDFs or image-based files without extractable text layers, necessitating manual page-by-page review.
\end{itemize}

\section{Importance of Legal NLP}
Legal Natural Language Processing (Legal NLP) is a specialized branch of artificial intelligence that focuses on analyzing, processing, and understanding legal documents. Unlike generic NLP systems trained on news articles or social media, Legal NLP addresses the unique characteristics of legal text:
\begin{itemize}
    \item \textbf{Syntactic Complexity:} Legal drafting features extremely long sentences, nested clauses, parenthetical explanations, and archaic vocabulary (often referred to as ``legalese''). Legal NLP tools are optimized to parse these complex syntax trees.
    \item \textbf{Semantic Precision:} In the legal domain, minor phrasing variations can alter a party's rights or obligations. Legal NLP systems use contextual representations to capture subtle semantic details that generic models overlook.
    \item \textbf{Contextual Disambiguation:} Words like ``consideration,'' ``performance,'' or ``execution'' have specific meanings in legal contexts that differ from their everyday usage. Legal NLP models resolve these ambiguities using surrounding context.
\end{itemize}

\section{Importance of Machine Learning}
Machine Learning (ML) acts as the driving engine behind automated document classification and search. By representing text numerically (vectorization), ML models learn patterns from training data to classify new, unseen texts. Instead of relying on manual rules or regular expressions, which fail when encountering variations in drafting, machine learning models learn representation boundaries. This enables automated classification of clauses with high generalization capabilities, reducing the manual review workload from days to seconds.

\section{Why Classical Machine Learning is Chosen instead of Deep Learning or RAG}
While deep learning models (such as BERT, RoBERTa, or LSTM) and Large Language Models (LLMs) used in Retrieval-Augmented Generation (RAG) have become popular in general NLP, classical machine learning remains the preferred choice for this project. The selection is justified by several practical and technical factors:
\begin{enumerate}
    \item \textbf{Interpretability and Transparency:} In legal applications, explainability is crucial. Lawyers must understand why a system classified a clause under a specific category. Classical ML algorithms like Logistic Regression and Naive Bayes provide clear coefficients and probabilities, and Random Forests output feature importances. In contrast, deep neural networks and LLMs function as black boxes, making error tracing highly difficult.
    \item \textbf{Computational Efficiency and Low Resource Footprint:} LLMs and deep learning models require expensive, high-end Graphical Processing Units (GPUs) for training and inference, which translates to high operational costs. Classical ML algorithms (such as Naive Bayes, Logistic Regression, Linear SVM, and Random Forests) are highly lightweight, enabling fast training and sub-millisecond inference on standard CPUs. This makes local desktop deployment viable and cost-effective.
    \item \textbf{Determinism and Prevention of Hallucination:} Generative AI models and RAG pipelines are prone to hallucinating facts, misinterpreting legal limits, or outputting text that does not exist in the source document. In legal compliance, a single hallucination can lead to significant liability. Classical ML models perform deterministic classification based strictly on the vocabulary present in the source text, eliminating the risk of hallucination.
    \item \textbf{Fit Within MCA Mini Project Scope:} Developing and optimizing deep learning models or setting up complex RAG orchestrations requires massive datasets, lengthy training times, and expensive APIs. Classical machine learning combined with Term Frequency-Inverse Document Frequency (TF-IDF) feature matrices is appropriate for the time constraints and academic scope of an MCA Mini Project. It demonstrates core machine learning principles—data preprocessing, feature extraction, classifier training, comparison, and evaluation—without unnecessary complexity.
\end{enumerate}

\section{Project Transition: From DocuMind AI to ClauseIQ}
The project was initially proposed under the title \textbf{DocuMind AI}. The initial concept was designed to build a general-purpose document intelligence system that could process multiple file types (PDFs, Word documents, images, and text files) across various domains (financial invoices, medical records, and legal files). 

However, during the initial feasibility study, this broad scope was found to be impractical within the time constraints of the MCA Mini Project. Trying to cover diverse document domains would result in shallow models and poor accuracy due to conflicting vocabularies. 

Consequently, the project scope was refined and focused exclusively on the legal domain under the new title \textbf{ClauseIQ}. ClauseIQ focuses on legal contracts, concentrating on automated clause extraction, 41-category classification, and intelligent similarity-based search. This narrowed focus allows for deep optimization of features and algorithms, leveraging domain-specific legal vocabularies to deliver high accuracy. A more ambitious deep-learning / RAG-based system (originally proposed as "DocuMind AI") was considered but was found to be too broad for this stage and has been deferred to a possible fourth-semester main project.

\section{Introduction to the CUAD Dataset}
ClauseIQ utilizes the \textbf{Contract Understanding Atticus Dataset (CUAD)}, developed by The Atticus Project. CUAD is a curated, expert-annotated corpus specifically designed for legal contract clause classification. It contains over 13,000 legal clauses extracted from 510 commercial contracts drawn from the U.S. Securities and Exchange Commission's public EDGAR database and expert-annotated by lawyers and law students associated with The Atticus Project. CUAD identifies 41 categories of clauses that are routinely reviewed in corporate transactions — such as Governing Law, Termination for Convenience, Non-Compete, and Cap on Liability — spanning more than 13,000 individually annotated clause instances. This makes it one of the largest, highest-quality, expert-labelled legal NLP resources publicly available, and a natural fit for a supervised clause-classification task.

\section{Two Key Objectives of ClauseIQ}
The ClauseIQ system is designed to achieve two primary functional objectives:
\begin{enumerate}
    \item \textbf{Contract Clause Classification:} Given the text of a clause (or a labelled span extracted from a contract), predict which of the 41 CUAD categories it belongs to, so that a reviewer can be automatically shown ``this looks like a Non-Compete clause'' instead of reading the whole document.
    \item \textbf{Intelligent Clause Search:} Given a natural-language query (e.g. ``find clauses that cap our liability''), retrieve the most relevant clauses across the corpus using TF-IDF vector representations and cosine similarity, without requiring an exact keyword match.
\end{enumerate}

Together, classification and search form a lightweight ``contract intelligence'' pipeline that mirrors, at a much smaller scale, the kind of tooling used in commercial contract-review products.

\section{Project Scope}
The scope of ClauseIQ includes:
\begin{itemize}
    \item Developing a preprocessing pipeline to clean raw contract texts.
    \item Implementing feature extraction using TF-IDF vectorization with unigrams and bigrams.
    \item Training and evaluating four classical machine learning classifiers: Multinomial Naive Bayes, Support Vector Machine (Linear SVM), Logistic Regression, and Random Forest.
    \item Developing a search module using Cosine Similarity to find relevant clauses.
    \item Providing a user-friendly interface to upload contracts, view classified clauses, and perform search queries.
\end{itemize}
This initial report documents the first phase of the project: dataset finalisation, supporting literature review, exploratory data analysis, preprocessing design, algorithm selection, project pipeline, and feasibility analysis. Model training, evaluation and deployment are addressed in the subsequent implementation sprints and reported in the final report.

\section{Project Objectives}
The core engineering objectives are:
\begin{itemize}
    \item To automate the tedious process of manual contract review.
    \item To achieve high classification accuracy and Macro F1-scores across all 41 categories in the CUAD dataset.
    \item To minimize query latency in the similarity search module.
    \item To build an explainable system that presents classification probabilities.
\end{itemize}

\section{Report Organization}
This initial report is organized into three main chapters. Chapter 1 provides the introduction, background, importance of Legal NLP, algorithmic choices, project transition, and objectives. Chapter 2 reviews the supporting literature, summarizing three relevant IEEE publications, followed by a comparative analysis and justification of the design decisions. Chapter 3 contains the System Analysis, detailing the CUAD dataset, preprocessing steps, algorithm workflows, pseudo-codes, and project pipeline.

\newpage

% ==========================================
% CHAPTER 2: SUPPORTING LITERATURE
% ==========================================
\chapter{SUPPORTING LITERATURE}
\label{ch:supporting_literature}

\section{Literature Review}

\subsection{Paper 1: LegalReasoner}
\noindent\textbf{Authors:} Wang X., Zhang X., Hoo V., Shao Z., Zhang X. (2024)\\
\textbf{Title:} ``LegalReasoner: A Multi-Stage Framework for Legal Judgment Prediction via Large Language Models and Knowledge Integration.''\\
\textbf{Journal:} \textit{IEEE Access}, vol. 12, pp. 166843--166854.\\
\textbf{DOI:} \url{https://doi.org/10.1109/ACCESS.2024.3496666}

\paragraph{Area of Work:}
The paper focuses on Legal Judgment Prediction (LJP). LJP is the task of predicting the outcome of legal cases, including applicable law articles, charges, and prison terms, based on the text of the case facts.

\paragraph{Dataset:}
The framework was evaluated using two prominent datasets: the European Court of Human Rights (ECHR) dataset and the Chinese AI and Law Challenge (CAIL2018) dataset, which is a benchmark dataset for judicial outcome prediction.

\paragraph{Methodology:}
The authors proposed a multi-stage framework consisting of four main stages: (i) legal knowledge infusion, where a large language model is pre-trained on legal literature using contrastive learning; (ii) case-law retrieval, using a graph neural network to surface relevant precedents and statutes; (iii) multi-hop reasoning over the retrieved material using a transformer architecture; and (iv) judgment synthesis.

\paragraph{Algorithms:}
The framework utilizes Large Language Models, Contrastive Learning, Graph Neural Networks, Transformer-based Hierarchical Attention, and Generative Adversarial Networks (GANs).

\paragraph{Results:}
The model showed a significant improvement over baseline models, achieving an average accuracy increase of 7.8\% on LJP tasks.

\paragraph{Advantages:}
The system integrates domain-specific legal knowledge, improving explainability and prediction accuracy for complex legal outcomes.

\paragraph{Future Scope:}
The authors suggest extending the framework to multi-lingual legal systems and incorporating real-time legal updates.

\begin{table}[h]
\centering
\small
\begin{tabular}{p{3.5cm}p{11.5cm}}
\toprule
\textbf{Parameter} & \textbf{Details} \\
\midrule
\textbf{Title} & LegalReasoner: A Multi-Stage Framework for Legal Judgment Prediction \\
\textbf{Area of Work} & Legal Judgment Prediction (LJP) \\
\textbf{Dataset} & ECHR and CAIL2018 datasets \\
\textbf{Methodology} & Four-stage framework: Knowledge Infusion, Case-Law Retrieval, Multi-Hop Reasoning, Synthesis \\
\textbf{Algorithms} & LLMs, Contrastive Learning, GNN, Hierarchical Attention, GANs \\
\textbf{Results} & 7.8\% average accuracy improvement over state-of-the-art baselines \\
\textbf{Advantages} & High accuracy, structured integration of legal knowledge, improved explainability \\
\textbf{Future Scope} & Support for multi-lingual legal texts, real-time case retrieval, dynamic statute updates \\
\bottomrule
\end{tabular}
\caption{Summary of Paper 1}
\label{tab:paper1_summary}
\end{table}

\newpage

\subsection{Paper 2: Exploring LLMs Applications in Law}
\noindent\textbf{Authors:} Siino M., Falco M., Croce D., Rosso P. (2025)\\
\textbf{Title:} ``Exploring LLMs Applications in Law: A Literature Review on Current Legal NLP Approaches.''\\
\textbf{Journal:} \textit{IEEE Access}, vol. 13, pp. 18253--18276.\\
\textbf{DOI:} \url{https://doi.org/10.1109/ACCESS.2025.3533217}

\paragraph{Area of Work:}
This is a systematic literature review (61 selected publications, 2017–2023) that maps how Transformer/LLM-based NLP has been applied across the legal domain, including legal document analysis, case prediction, legal search, regulatory compliance and contract review.

\paragraph{Dataset:}
The paper reviews 61 key publications published between 2017 and 2023, analyzing benchmarks like CUAD, ContractNLI and other legal NLP corpora.

\paragraph{Methodology:}
The authors conducted a systematic mapping study to categorize legal NLP research into key application areas, identifying methodological gaps, evaluation protocols, and trends in model sizes and fine-tuning strategies.

\paragraph{Algorithms:}
The review covers various transformer-based architectures including BERT, RoBERTa, legal-specific models like Legal-BERT, CaseBERT, T5, and GPT-based generative models.

\paragraph{Results:}
The paper details consistent improvements across legal NLP tasks like document summarization, QA, and clause extraction. Fine-tuning models on domain-specific data (e.g., Legal-BERT) is shown to consistently outperform generic models of comparable size.

\paragraph{Advantages:}
The review provides a structured overview of legal NLP, highlighting open challenges in model hallucinations, privacy, and computational costs.

\paragraph{Future Scope:}
The authors highlight the need for research into domain-specific prompt engineering, mitigating hallucinations, integrating privacy-preserving architectures (like federated learning), and developing hybrid models that combine LLMs with symbolic legal reasoning.

\begin{table}[h]
\centering
\small
\begin{tabular}{p{3.5cm}p{11.5cm}}
\toprule
\textbf{Parameter} & \textbf{Details} \\
\midrule
\textbf{Title} & Exploring LLMs Applications in Law: A Literature Review \\
\textbf{Area of Work} & Systematic review of Legal NLP and Transformer-based models \\
\textbf{Dataset} & Survey of 61 core publications (utilizing CUAD, ContractNLI, etc.) \\
\textbf{Methodology} & PRISMA-style systematic review, taxonomy development, and gap analysis \\
\textbf{Algorithms} & BERT, RoBERTa, Legal-BERT, CaseBERT, T5, GPT-based LLMs \\
\textbf{Results} & Documents performance trends; confirms Legal-BERT outperforming generic baselines \\
\textbf{Advantages} & Highlights challenges like hallucinations, privacy concerns, and computational costs \\
\textbf{Future Scope} & Focus on domain-specific prompting, privacy preservation, and hybrid symbolic models \\
\bottomrule
\end{tabular}
\caption{Summary of Paper 2}
\label{tab:paper2_summary}
\end{table}

\newpage

\subsection{Paper 3: MCNN-LSTM}
\noindent\textbf{Authors:} Hasib K. M., Azam S., Karim A., Al Marouf A., Shamrat F. M. J., Montaha S., Yeo K. C., Jonkman M., Alhajj R., Rokne J. G. (2023)\\
\textbf{Title:} ``MCNN-LSTM: Combining CNN and LSTM to Classify Multi-Class Text in Imbalanced News Data.''\\
\textbf{Journal:} \textit{IEEE Access}, vol. 11, pp. 93048--93063.\\
\textbf{DOI:} \url{https://doi.org/10.1109/ACCESS.2023.3309697}

\paragraph{Area of Work:}
The paper focuses on multi-class text classification in imbalanced datasets, which is highly relevant to legal clause datasets like CUAD that exhibit severe category imbalance.

\paragraph{Dataset:}
Evaluated on imbalanced news datasets (including BBC News and Reuters) containing articles categorized into multiple news segments.

\paragraph{Methodology:}
The authors proposed a hybrid deep learning model combining a Multi-channel Convolutional Neural Network (MCNN) and a Long Short-Term Memory (LSTM) network. The model extracts local n-gram features and sequential long-range dependencies, utilizing class weights and Tomek-Link sampling to handle dataset imbalance.

\paragraph{Algorithms:}
The model uses Multi-channel CNN, LSTM, word embeddings (Word2Vec, GloVe), and a Softmax classifier.

\paragraph{Results:}
The model achieved an F1-score of 98\% and an accuracy of 99.71\%, outperforming standard baseline models on imbalanced text datasets.

\paragraph{Advantages:}
Effectively captures both local syntax patterns and long-range contextual relationships, while maintaining high performance on minority classes.

\paragraph{Future Scope:}
The authors suggest testing the model on other domains (such as legal text) and exploring Transformer layers instead of LSTM.

\begin{table}[h]
\centering
\small
\begin{tabular}{p{3.5cm}p{11.5cm}}
\toprule
\textbf{Parameter} & \textbf{Details} \\
\midrule
\textbf{Title} & MCNN-LSTM: Combining CNN and LSTM to Classify Multi-Class Text \\
\textbf{Area of Work} & Multi-class text classification under severe class imbalance \\
\textbf{Dataset} & Imbalanced news datasets (BBC News, Reuters) \\
\textbf{Methodology} & Parallel CNN filters for local patterns, LSTM for sequential context, class-balancing \\
\textbf{Algorithms} & MCNN, LSTM, Word2Vec, GloVe, Softmax \\
\textbf{Results} & 99.71\% classification accuracy and 98\% F1-score \\
\textbf{Advantages} & Captures local and global text representations; handles imbalanced classes \\
\textbf{Future Scope} & Evaluation on non-news datasets, replacement of LSTM with Transformer layers \\
\bottomrule
\end{tabular}
\caption{Summary of Paper 3}
\label{tab:paper3_summary}
\end{table}

\newpage

\section{Literature Review Summary}

\begin{table}[h]
\centering
\scriptsize
\begin{tabular}{p{2cm}p{2.3cm}p{2.3cm}p{2.3cm}p{2.3cm}p{2.3cm}p{2.3cm}}
\toprule
\textbf{Title} & \textbf{Area} & \textbf{Dataset} & \textbf{Methodology} & \textbf{Algorithms} & \textbf{Advantages} & \textbf{Limitations} \\
\midrule
\textbf{LegalReasoner} (2024) & Legal Judgment Prediction & ECHR, CAIL2018 & Multi-stage knowledge infusion and graph-based retrieval & LLMs, GNNs, GANs, Hierarchical Attention & High accuracy; incorporates legal context & Expensive GPUs; high latency \\
\textbf{LLMs in Law} (2025) & Systematic mapping of Legal NLP & Survey of 61 papers & Systematic review of Transformer applications & BERT, RoBERTa, Legal-BERT, GPT & Maps research trends and gaps & Review-only; no new model implementation \\
\textbf{MCNN-LSTM} (2023) & Text classification in imbalanced datasets & News datasets (BBC, Reuters) & Multi-channel feature extraction and class-weighting & CNN, LSTM, Word2Vec, Softmax & High F1-score; robust local/global features & Computationally heavy; hard to interpret \\
\bottomrule
\end{tabular}
\caption{Literature Review Comparison Table}
\label{tab:literature_summary}
\end{table}

\section{Findings and Proposals}
The reviewed literature converges on three points that directly shape ClauseIQ's design. First, contract clause classification and semantic clause search are recognised, well-motivated sub-tasks within legal NLP (Siino et al., 2025), and CUAD is an established benchmark for them. Second, pairing a classification component with a retrieval/search component — as LegalReasoner does with GNN-based case-law retrieval alongside its reasoning stage (Wang et al., 2024) — is a proven pattern for structured legal text systems; ClauseIQ mirrors this at a smaller scale by pairing a TF-IDF + classical-ML classifier with a TF-IDF + cosine-similarity search module. Third, and most practically, Hasib et al. (2023) show that on multi-class, heavily imbalanced text classification, classical ML models (Logistic Regression, Naive Bayes, SVM) remain reasonably competitive once imbalance-aware evaluation and training strategies are applied, even though a deep hybrid model ultimately scores highest. Since CUAD is itself severely imbalanced across its 41 categories, this finding directly justifies ClauseIQ's decision to build a classical ML baseline first — using TF-IDF features with Naive Bayes, SVM, Logistic Regression and Random Forest, evaluated with macro-averaged F1 and per-class metrics — rather than committing immediately to a deep-learning or transformer-based pipeline.

Based on these findings, the proposed direction for ClauseIQ is: (1) build a clean, deduplicated clause-level dataset from CUAD's 41 categories; (2) represent clause text using TF-IDF (unigrams and bigrams); (3) train and compare Multinomial Naive Bayes, Linear SVM, Logistic Regression and Random Forest classifiers, using class-weighting and macro-F1 to account for the pronounced class imbalance documented in Section 3; (4) build a companion intelligent-search module that ranks clauses by cosine similarity to a free-text query over the same TF-IDF space; and (5) reserve deep-learning / transformer-based extensions (e.g., Legal-BERT, as flagged in the literature) for a possible future main project, consistent with the scope constraint that led the mentor to reject the earlier, more ambitious DocuMind AI proposal.

\newpage

% ==========================================
% CHAPTER 3: SYSTEM ANALYSIS
% ==========================================
\chapter{SYSTEM ANALYSIS}
\label{ch:system_analysis}

All figures, statistics and code outputs in this section were produced by actually loading and analysing the real CUAD v1 dataset (obtained from the official Atticus Project GitHub repository, \url{https://github.com/TheAtticusProject/cuad}) rather than a simulated or placeholder dataset.

\section{Analysis of the Dataset}

\subsection{About the Dataset}
The Contract Understanding Atticus Dataset (CUAD) v1 is sourced from the Atticus Project and distributed at \url{https://github.com/TheAtticusProject/cuad} (and mirrored on Hugging Face Datasets). It contains 510 real commercial contracts collected from the U.S. Securities and Exchange Commission's public EDGAR filing system, spanning 25 contract types (Distributor Agreements, Franchise Agreements, Licence Agreements, Strategic Alliance Agreements, etc.). Each contract has been manually annotated by law-student and attorney annotators against 41 clause categories that lawyers routinely review during due diligence — for example Governing Law, Anti-Assignment, Cap On Liability, Non-Compete, and Ip Ownership Assignment. In its raw distributed form, CUAD is released as a SQuAD-style question-answering JSON file (\texttt{CUADv1.json}), where each of the 41 categories is phrased as a ``question'' (e.g. ``Highlight the parts of this contract related to ‘Non-Compete’...'') and the answer is one or more text spans (the actual clause text) located inside the contract.

For ClauseIQ's classification and search use case, this QA-style format was converted into a flat, clause-level table: each row is one annotated clause span, together with the contract it came from and its category label. This conversion was performed directly on the official \texttt{CUADv1.json} file using a Python script (\texttt{json} + \texttt{pandas}), producing the working dataset \texttt{clauseiq\_dataset.csv} used for all analysis in this report.
\begin{itemize}
    \item \textbf{Dataset source:} \url{https://github.com/TheAtticusProject/cuad}
    \item \textbf{Raw file used:} \texttt{CUADv1.json} (SQuAD-style question-answering format)
    \item \textbf{Contracts:} 510 real commercial contracts sourced from SEC EDGAR
    \item \textbf{Clause categories (classes):} 41
    \item \textbf{Annotated clause spans after flattening and de-duplication:} 13,823
    \item \textbf{License:} Creative Commons Attribution 4.0 (CC BY 4.0)
\end{itemize}

\subsection{Explore the Dataset}
The flattened clause-level dataset was loaded with pandas and inspected using the same exploratory workflow used for structured tabular data: shape, data types, first records, missing values and summary statistics.

\paragraph{a) Shape and data types}
The database contains 13,823 rows and 5 columns:
\begin{table}[h]
\centering
\small
\begin{tabular}{p{4cm}p{11cm}}
\toprule
\textbf{Property} & \textbf{Details} \\
\midrule
\textbf{Shape} & 13,823 rows, 5 columns \\
\textbf{Columns} & \texttt{contract}, \texttt{category}, \texttt{clause\_text}, \texttt{char\_len}, \texttt{word\_count} \\
\textbf{Datatypes} & \texttt{contract}: object, \texttt{category}: object, \texttt{clause\_text}: object, \texttt{char\_len}: int64, \texttt{word\_count}: int64 \\
\bottomrule
\end{tabular}
\caption{Dataset Schema and Shape summary}
\label{tab:dataset_shape}
\end{table}

\begin{figure}[h]
\centering
\framebox[0.85\textwidth]{\vbox to 1cm{\vfill\centering \texttt{SHAPE: (13823, 5)}\\ \texttt{contract: object, category: object, clause\_text: object, char\_len: int64, word\_count: int64}\vfill}}
\caption{Shape and data types of the flattened clause dataset}
\label{fig:shape_datatypes}
\end{figure}

\paragraph{b) First five records}
\begin{figure}[h]
\centering
\framebox[0.85\textwidth]{\vbox to 2cm{\vfill\centering \scriptsize
\texttt{0 LIMEENERGYCO\_09\_09\_1999-EX-10-DISTRIBUTOR AGR... Document Name DISTRIBUTOR AGREEMENT 21 2} \\
\texttt{1 LIMEENERGYCO\_09\_09\_1999-EX-10-DISTRIBUTOR AGR... Parties Distributor 11 1} \\
\texttt{2 LIMEENERGYCO\_09\_09\_1999-EX-10-DISTRIBUTOR AGR... Parties Electric City Corp. 19 3} \\
\texttt{3 LIMEENERGYCO\_09\_09\_1999-EX-10-DISTRIBUTOR AGR... Parties Electric City of Illinois L.L.C. 32 5} \\
\texttt{4 LIMEENERGYCO\_09\_09\_1999-EX-10-DISTRIBUTOR AGR... Parties Company 7 1}
\vfill}}
\caption{First five rows of the clause-level dataset}
\label{fig:first_five}
\end{figure}

\paragraph{c) Missing values and duplicates}
Because every row was constructed only from non-empty answer spans in the source JSON, the flattened dataset contains 0 missing values and 0 exact duplicate rows.
\begin{figure}[h]
\centering
\framebox[0.85\textwidth]{\vbox to 1.5cm{\vfill\centering \texttt{MISSING VALUES: contract: 0, category: 0, clause\_text: 0, char\_len: 0, word\_count: 0} \\ \texttt{DUPLICATE ROWS: 0}\vfill}}
\caption{Missing value and duplicate-row check}
\label{fig:missing_duplicates}
\end{figure}

\paragraph{d) Summary statistics}
\begin{table}[h]
\centering
\small
\begin{tabular}{p{3.5cm}p{3cm}p{3cm}p{3cm}p{2.5cm}}
\toprule
\textbf{Metric} & \textbf{Mean} & \textbf{Std Dev} & \textbf{Min} & \textbf{Max} \\
\midrule
\textbf{Char Length} & 262.01 chars & 285.15 chars & 1.0 char & 3,884 chars \\
\textbf{Word Count} & 40.67 words & 43.44 words & 1.0 word & 479 words \\
\bottomrule
\end{tabular}
\caption{Summary Statistics of Clause Lengths}
\label{tab:summary_stats}
\end{table}

\begin{figure}[h]
\centering
\framebox[0.85\textwidth]{\vbox to 1.5cm{\vfill\centering \scriptsize \texttt{count: 13823, mean(word): 40.68, std(word): 43.45, min(word): 1, 50\%(word): 31, max(word): 479}\vfill}}
\caption{Summary statistics of clause length (characters and words)}
\label{fig:summary_stats}
\end{figure}

Clause length is right-skewed: the median clause is 31 words long, but a long tail of clauses (up to 479 words) corresponds to lengthy definitional or indemnification-style provisions. This motivates using TF-IDF (which is length-normalised) rather than raw word counts as the feature representation.

\paragraph{e) Clauses per contract and per category}
On average each contract contributes about 27.1 labelled clauses, but this varies widely (2 to 97), reflecting genuine differences in contract length and complexity across the 25 contract types in CUAD.
\begin{figure}[h]
\centering
\framebox[0.85\textwidth]{\vbox to 1.5cm{\vfill\centering \scriptsize \texttt{count: 510, mean: 27.1, std: 18.12, min: 2, 50\%: 22, max: 97}\vfill}}
\caption{Distribution of number of annotated clauses per contract}
\label{fig:clauses_per_contract}
\end{figure}

The category distribution is highly uneven: ``Parties'' alone accounts for 2,554 of the 13,823 clauses (about 18\%), followed by ``License Grant'', ``Cap On Liability'', ``Anti-Assignment'' and ``Audit Rights'', each in the 600–800 range. This class imbalance, discussed further in Section \ref{subsec:class_analysis}, is a central design consideration for the modelling stage.

\begin{figure}[h]
\centering
\framebox[0.85\textwidth]{\vbox to 2.5cm{\vfill\centering [Top 15 Categories: Parties (2554), License Grant (752), Cap On Liability (680), ...]\vfill}}
\caption{Top 15 most frequent clause categories in CUAD}
\label{fig:top_15_categories}
\end{figure}

We also outline the following visual placeholders:
\begin{figure}[h]
\centering
\framebox[0.85\textwidth]{\vbox to 2cm{\vfill\centering [Histogram showing skewness of word count distribution]\vfill}}
\caption{Distribution of clause length (word count)}
\label{fig:length_dist}
\end{figure}

\begin{figure}[h]
\centering
\framebox[0.85\textwidth]{\vbox to 2cm{\vfill\centering [Boxplot representing clause distributions over 510 contracts]\vfill}}
\caption{Number of annotated clauses per contract}
\label{fig:clauses_per_contract_boxplot}
\end{figure}

\newpage

\section{Data Preprocessing}

\subsection{Data Cleaning}
\label{subsec:data_cleaning}
Two cleaning steps were applied to the flattened clause dataset before feature extraction:
\begin{enumerate}
    \item \textbf{Text normalisation:} collapsing repeated whitespace/newlines (a common OCR/formatting artefact in the source contracts) and stripping characters outside a legal-text-safe set (letters, digits, and common punctuation such as periods, commas, parentheses, \%, \$, and hyphens).
    \item \textbf{Removal of extremely short spans:} clause spans under 3 tokens (e.g. a lone party name or date fragment picked up as a partial answer span) carry little topical signal for a 41-way classifier and were removed.
\end{enumerate}

Before cleaning: 13,823 rows. Removed 1,567 clauses with < 3 tokens (11.34\%). Cleaned dataset shape is (12,256, 6).

\begin{figure}[h]
\centering
\framebox[0.85\textwidth]{\vbox to 1.5cm{\vfill\centering \texttt{Before cleaning: (13823, 5)} \\ \texttt{Removed 1567 clauses with < 3 tokens (11.34\%)} \\ \texttt{Cleaned dataset shape: (12256, 6)}\vfill}}
\caption{Result of data cleaning and short-span removal}
\label{fig:cleaning_results}
\end{figure}

\begin{lstlisting}[language=Python, caption={Text Cleaning and Normalization Code Snippet}]
import re

def clean_text(t):
    t = re.sub(r'\s+', ' ', str(t))                  # collapse whitespace
    t = re.sub(r"[^A-Za-z0-9.,;:()%$\-\s]", ' ', t)  # strip odd artefacts
    return t.strip()

df['clean_text'] = df['clause_text'].apply(clean_text)
df_clean = df[df['clean_text'].str.split().str.len() >= 3]
\end{lstlisting}

\subsection{Analysis of Feature Variables}
\label{subsec:feature_analysis}
Clause text is converted into numerical features using Term Frequency–Inverse Document Frequency (TF-IDF) with unigrams and bigrams, a vocabulary capped at 5,000 terms, English stop-word removal, and a minimum document frequency of 2 (to discard one-off tokens such as OCR noise or party names).

\begin{lstlisting}[language=Python, caption={TF-IDF Feature Extraction Code Snippet}]
from sklearn.feature_extraction.text import TfidfVectorizer

tfidf = TfidfVectorizer(max_features=5000, stop_words='english',
                        ngram_range=(1,2), min_df=2)
X = tfidf.fit_transform(df_clean['clean_text'])
print('TF-IDF matrix shape:', X.shape)
print('Sparsity: {:.4f}%'.format(
    100 * (1 - X.nnz / (X.shape[0]*X.shape[1]))))
\end{lstlisting}

\begin{figure}[h]
\centering
\framebox[0.85\textwidth]{\vbox to 1.8cm{\vfill\centering \texttt{TF-IDF matrix shape: (12256, 5000)} \\ \texttt{Sparsity: 99.5187\%} \\ \scriptsize \texttt{Top terms: agreement(449.05), shall(394.78), party(353.97), term(211.92), company(203.05)}\vfill}}
\caption{TF-IDF feature matrix statistics and top terms}
\label{fig:tfidf_stats}
\end{figure}

\begin{figure}[h]
\centering
\framebox[0.85\textwidth]{\vbox to 2cm{\vfill\centering [Bar chart of top 20 terms: agreement, shall, party, term, company, ...]\vfill}}
\caption{Top 20 TF-IDF terms across the clause corpus}
\label{fig:top_terms_bar}
\end{figure}

As expected for legal text, the highest-weighted terms are boilerplate contract vocabulary (``agreement'', ``shall'', ``party'', ``section''). These generic terms are shared across almost every category, so the classifiers must rely on more category-specific bigrams (e.g. ``agreement shall'', ``prior written'') and rarer discriminating terms to separate the 41 classes — this is one reason bigrams were included in the TF-IDF configuration rather than unigrams alone.

\subsection{Analysis of Class Variables}
\label{subsec:class_analysis}
\begin{figure}[h]
\centering
\framebox[0.85\textwidth]{\vbox to 2cm{\vfill\centering [Histogram of categories grouped by sample count bins (e.g., 0-100, 101-200)]\vfill}}
\caption{Class imbalance: number of categories by sample-count bucket}
\label{fig:class_imbalance_bins}
\end{figure}

The 41 clause categories are strongly imbalanced. The largest class (``Parties'', 2,554 samples) has roughly 95 times more examples than the smallest class (``Price Restrictions'', 27 samples). Four categories fall below 40 samples after cleaning:
\begin{figure}[h]
\centering
\framebox[0.85\textwidth]{\vbox to 1.5cm{\vfill\centering \small \texttt{Third Party Beneficiary: 39} \\ \texttt{Most Favored Nation: 38} \\ \texttt{Unlimited/All-You-Can-Eat-License: 32} \\ \texttt{Price Restrictions: 27}\vfill}}
\caption{Rare categories identified for class-imbalance handling}
\label{fig:rare_categories}
\end{figure}

Following the imbalance-handling strategy validated by Hasib et al. (2023) in Section \ref{ch:supporting_literature}, the modelling stage will address this imbalance using class-weighted loss (\texttt{class\_weight='balanced'} in scikit-learn), stratified train/test splitting so every category is represented in both splits, and macro-averaged F1-score (rather than plain accuracy) as the primary evaluation metric, since accuracy alone would be dominated by the majority ``Parties'' class.

\section{Data Visualization}
The following visualisations summarise the shape of the CUAD-derived clause dataset used for ClauseIQ.
\begin{figure}[h]
\centering
\framebox[0.85\textwidth]{\vbox to 2.2cm{\vfill\centering [Bar chart representing the full distribution of samples across all 41 categories]\vfill}}
\caption{Full distribution of all 41 clause categories}
\label{fig:full_category_dist}
\end{figure}

\begin{figure}[h]
\centering
\framebox[0.85\textwidth]{\vbox to 2.2cm{\vfill\centering [Boxplot of word lengths for the top 10 most frequent categories]\vfill}}
\caption{Clause length by category (top 10 most frequent categories)}
\label{fig:length_by_category}
\end{figure}

Clause length also varies meaningfully by category — e.g. ``Audit Rights'' and ``Insurance'' clauses tend to be longer and more detailed than short factual spans such as ``Agreement Date'' — which is useful contextual information even though the primary features used for classification are TF-IDF weighted terms rather than length.

\begin{figure}[h]
\centering
\framebox[0.85\textwidth]{\vbox to 2.2cm{\vfill\centering [Train/Test Split ratio diagram: 80\% training split, 20\% validation split, stratified]\vfill}}
\caption{Planned stratified train/test split of the clause dataset}
\label{fig:train_test_split_diagram}
\end{figure}

A stratified 80:20 train/test split (\texttt{random\_state = 42}) is planned for the modelling stage so that every one of the 41 categories — including the rarest ones — is represented proportionally in both the training and test sets.

\newpage

\section{Analysis of Architecture and Algorithms}

\subsection{Detailed Study Report of the Proposed Algorithm(s)}
ClauseIQ combines a TF-IDF feature representation with four classical supervised classifiers for clause categorisation, and a TF-IDF + cosine-similarity module for intelligent search over the same feature space.

\subsubsection{TF-IDF Vectorisation}
Term Frequency–Inverse Document Frequency converts each clause into a sparse numeric vector where each dimension corresponds to a term (unigram or bigram); a term's weight is high if it occurs often in that clause but rarely across the corpus, which surfaces category-discriminating vocabulary and down-weights generic boilerplate. TF-IDF is well suited to short, boilerplate-heavy legal clauses because it does not require the large labelled corpora that transformer embeddings need, and it produces an interpretable feature space.

Mathematical formulations:
\begin{equation}
    TF(t, d) = \frac{f_{t,d}}{\sum_{t' \in d} f_{t', d}}
\end{equation}
\begin{equation}
    IDF(t, D) = \log \left( \frac{1 + N}{1 + DF(t)} \right) + 1
\end{equation}
\begin{equation}
    TF\text{-}IDF(t, d, D) = TF(t, d) \times IDF(t, D)
\end{equation}

\subsubsection{Multinomial Naive Bayes}
A probabilistic classifier that assumes feature (term) independence given the class. It is extremely fast to train, performs well on sparse high-dimensional TF-IDF matrices, and serves as a strong, simple baseline for multi-class text classification.
\begin{equation}
    P(c|d) \propto P(c) \prod_{i=1}^{n} P(t_i|c)
\end{equation}
\begin{equation}
    P(t_i|c) = \frac{count(t_i, c) + \alpha}{\sum_{t \in V} count(t, c) + \alpha |V|}
\end{equation}

\subsubsection{Support Vector Machine (Linear Kernel)}
SVMs find the maximum-margin hyperplane separating classes and are known to perform particularly well on high-dimensional, sparse text data such as TF-IDF vectors, where the number of features (up to 5,000) is much larger than what many other classifiers can handle efficiently. A linear kernel and one-vs-rest scheme are used to extend the binary SVM to the 41-class problem.
\begin{equation}
    \min_{w, b, \xi} \frac{1}{2} \|w\|^2 + C \sum_{i=1}^{n} \xi_i
\end{equation}
subject to $y_i (w^T x_i + b) \ge 1 - \xi_i, \quad \xi_i \ge 0$.

\subsubsection{Logistic Regression}
A linear probabilistic classifier trained with multinomial (softmax) loss, offering class probabilities (useful for ranking candidate categories, not just a single hard prediction) and good interpretability through per-class term coefficients, at low computational cost.
\begin{equation}
    P(y=c|x) = \frac{e^{w_c^T x + b_c}}{\sum_{j=1}^{C} e^{w_j^T x + b_j}}
\end{equation}

\subsubsection{Random Forest}
An ensemble of decision trees trained on bootstrap samples and random feature subsets, whose predictions are aggregated by majority vote. Random Forest can capture non-linear interactions between TF-IDF features and is naturally robust to noisy or irrelevant terms.
\begin{equation}
    \text{Gini}(D) = 1 - \sum_{i=1}^{C} p_i^2
\end{equation}

\subsubsection{Cosine Similarity (Intelligent Search)}
For the search module, both the corpus of clauses and an incoming natural-language query are projected into the same TF-IDF space, and clauses are ranked by the cosine of the angle between their vector and the query vector. Cosine similarity is scale-invariant (it is unaffected by clause length) and is the standard similarity measure for sparse bag-of-words / TF-IDF retrieval.
\begin{equation}
    \text{Similarity}(A, B) = \frac{A \cdot B}{\|A\| \|B\|} = \frac{\sum_{i=1}^n A_i B_i}{\sqrt{\sum_{i=1}^n A_i^2} \sqrt{\sum_{i=1}^n B_i^2}}
\end{equation}

\subsubsection{Performance Evaluation Metrics}
Model performance is validated quantitatively using:
\begin{itemize}
    \item \textbf{Precision:} $TP / (TP + FP)$
    \item \textbf{Recall:} $TP / (TP + FN)$
    \item \textbf{Accuracy:} $(TP + TN) / (TP + TN + FP + FN)$
    \item \textbf{F1-score:} $2 \times (\text{Precision} \times \text{Recall}) / (\text{Precision} + \text{Recall})$
    \item \textbf{Macro F1-score:} $\frac{1}{C} \sum_{c=1}^{C} F1_c$ (Primary indicator due to imbalance)
\end{itemize}

\subsection{Workflows and Pipeline Details}

\begin{figure}[h]
\centering
\framebox[0.85\textwidth]{\vbox to 3cm{\vfill\centering [Activity Diagram showing the process: Text Input -> Cleaning -> TF-IDF -> Classifier -> Class Prediction / Search Ranking]\vfill}}
\caption{Activity diagram: TF-IDF + classical ML clause classification and search pipeline}
\label{fig:activity_diagram}
\end{figure}

\begin{lstlisting}[language=Python, caption={Pipeline Training and Inference Pseudocode}]
# TF-IDF + Classifier Training and Similarity Searching

def train_clause_classifier(D, model_type):
    # D: Clause dataset with text X and labels Y
    X_clean = clean_text(D.text)
    tfidf = fit_tfidf(X_clean, max_features=5000, ngram_range=(1,2))
    X_train, X_test, y_train, y_test = stratified_split(tfidf, D.category, test_size=0.2)
    model = init_model(model_type, class_weight='balanced')   # NB / SVM / LR / RF
    model.fit(X_train, y_train)
    y_pred = model.predict(X_test)
    metrics = evaluate(y_test, y_pred)   # macro-F1, per-class F1, accuracy
    return model, tfidf, metrics

def search_clauses(query, corpus_vectors, tfidf, top_k):
    q_vec = tfidf.transform([query])
    scores = cosine_similarity(q_vec, corpus_vectors)
    return top_k_ranked(scores, top_k)
\end{lstlisting}

\newpage

\section{Project Pipeline}

\subsection{Project Pipeline Diagram with Explanation}
\begin{figure}[h]
\centering
\framebox[0.85\textwidth]{\vbox to 3.5cm{\vfill\centering [Pipeline Diagram: Raw CUAD JSON -> Flattened DataFrame -> Cleaning -> Split -> Model Compare -> Expose API]\vfill}}
\caption{ClauseIQ project pipeline}
\label{fig:project_pipeline}
\end{figure}

\paragraph{Data}
The pipeline begins with the CUAD dataset (510 contracts, 41 clause categories), which is flattened from its SQuAD-style JSON format into a clause-level table as described in Section \ref{ch:system_analysis}. Exploratory data analysis (class distribution, length distribution, missing-value and duplicate checks) is performed to understand the corpus, followed by data cleaning and TF-IDF-based preprocessing as described in Section \ref{subsec:data_cleaning}.

\paragraph{Model Building}
The cleaned, vectorised dataset is split into stratified training (80\%) and testing (20\%) subsets. Four classical ML models — Multinomial Naive Bayes, Linear SVM, Logistic Regression and Random Forest — are trained on the TF-IDF features with class-weighting to address the imbalance documented in Section \ref{subsec:class_analysis}, and compared using macro-F1, per-class F1 and accuracy to select the best-performing classifier.

\paragraph{Intelligent Search \& Deployment}
Independently of classification, a TF-IDF index of the full clause corpus is built to support cosine-similarity-based intelligent search: a free-text query is vectorised in the same TF-IDF space and the most similar clauses are ranked and returned. The selected classifier and the search index are planned to be persisted (via joblib) and exposed through a simple web/API interface so a user can classify a new clause or search the corpus for clauses matching a natural-language query.

\section{Feasibility Analysis}
Feasibility has been assessed across technical, economic and operational dimensions, consistent with the scope-controlled, classical-ML approach agreed with the project guide.

\subsection{Technical Feasibility}
\begin{itemize}
    \item \textbf{Data processing:} CUAD is a public, well-documented, CC-BY-4.0-licensed dataset; it was downloaded directly from the official GitHub repository and successfully flattened into a clean, tabular clause dataset (13,823 rows) using standard Python libraries (\texttt{pandas}, \texttt{json}, \texttt{re}), confirming that data acquisition and preprocessing are fully achievable within the project's technical resources.
    \item \textbf{Machine learning:} TF-IDF vectorisation and the four chosen classifiers (Naive Bayes, SVM, Logistic Regression, Random Forest) are all implemented in \texttt{scikit-learn}, a mature, well-documented, free library; no custom low-level ML code or GPU hardware is required.
    \item \textbf{Search module:} cosine similarity over TF-IDF vectors uses the same \texttt{scikit-learn} feature space already built for classification, so no additional infrastructure (e.g. a dedicated vector database) is required at this project's scale (13,823 clauses).
    \item \textbf{Deployment:} the trained classifier and TF-IDF index can be serialized with \texttt{joblib} and served through a lightweight Python web framework (e.g. FastAPI), consistent with tooling already used elsewhere in the department's mini projects.
\end{itemize}

\subsection{Economic Feasibility}
All software used — Python, pandas, scikit-learn, matplotlib/seaborn, Flask/FastAPI, joblib — is free and open-source; the CUAD dataset itself is free and openly licensed (CC BY 4.0).
No paid annotation, cloud GPU, or third-party API costs are required for the classical ML approach adopted here, keeping the project's running cost close to zero beyond standard computer hardware and development time.
The resulting tool has clear practical value — faster clause triage for legal/compliance teams, students, and paralegals reviewing standard commercial contracts — which justifies the (minimal) development effort.

\subsection{Operational Feasibility}
End users (e.g. a paralegal or contract reviewer) only need to paste clause text or a search query into a simple web interface; no legal-technology or machine-learning expertise is required to operate the system.
Because the underlying models are lightweight classical classifiers rather than large neural networks, inference is fast enough for interactive, single-clause or single-query use on ordinary hardware, without needing a GPU or specialised serving infrastructure.
The system's scope is intentionally bounded to the 41 CUAD categories and to classical ML, which keeps the operational surface small and maintainable within a single-semester mini project timeline, while leaving a clear, well-justified extension path (transformer / RAG-based deep learning) for a future main project.

\newpage

% ==========================================
% REFERENCES
% ==========================================
\begin{thebibliography}{99}

\bibitem{legalreasoner}
Wang X., Zhang X., Hoo V., Shao Z., Zhang X.,
``LegalReasoner: A Multi-Stage Framework for Legal Judgment Prediction via Large Language Models and Knowledge Integration,''
\textit{IEEE Access}, vol. 12, pp. 166843--166854, 2024.
\url{https://doi.org/10.1109/ACCESS.2024.3496666}

\bibitem{siino2025}
Siino M., Falco M., Croce D., Rosso P.,
``Exploring LLMs Applications in Law: A Literature Review on Current Legal NLP Approaches,''
\textit{IEEE Access}, vol. 13, pp. 18253--18276, 2025.
\url{https://doi.org/10.1109/ACCESS.2025.3533217}

\bibitem{hasib2023}
Hasib K. M., Azam S., Karim A., Al Marouf A., Shamrat F. M. J., Montaha S., Yeo K. C., Jonkman M., Alhajj R., Rokne J. G.,
``MCNN-LSTM: Combining CNN and LSTM to Classify Multi-Class Text in Imbalanced News Data,''
\textit{IEEE Access}, vol. 11, pp. 93048--93063, 2023.
\url{https://doi.org/10.1109/ACCESS.2023.3309697}

\bibitem{cuad}
D. Hendrycks, C. Burns, A. Chen, and S. Ball,
``CUAD: An Expert-Annotated NLP Dataset for Legal Contract Review,''
\textit{Proceedings of the Neural Information Processing Systems (NeurIPS)}, 2021.
\url{https://arxiv.org/abs/2103.06268}

\bibitem{scikit-learn}
F. Pedregosa et al.,
``Scikit-learn: Machine Learning in Python,''
\textit{Journal of Machine Learning Research}, vol. 12, pp. 2825--2830, 2011.

\bibitem{spacy}
M. Honnibal and I. Montani,
``spaCy: Industrial-strength Natural Language Processing in Python,''
\textit{Zenodo}, 2017.

\end{thebibliography}

\end{document}
